\section{Two class inputs for the pure-cubic quarter residues}
\label{sec:pure-cubic-quartic-class}

\paragraph{Status and exact scope.}
The finite unit, saturation and interval computations below have
separate author and independent review records.  Their combination
proves the displayed class-sector statements only after those records
are complete.  The provisional PARI class numbers are selectors for
integer isolation; neither a complete unit basis nor a provisional
class group is used as an input.  There is no original-coefficient
strict or common-window conclusion in this class theorem.

Put $p=269$, $K=\mathbb Q(\omega)$, $t=p^{1/3}>0$ and
$i^2=-1$, with $i\bmod p=82$.
Let $\psi=\omega_p^{67}$ be the primitive odd quartic character
modulo $p$, with $\psi(2)=-i$.
Let $\epsilon_K=\chi_{-3}$ and $\tau=\psi^2$.
Use the actual Gaussian-period field and its real twist:
\[
 \begin{aligned}
 F_-&=\mathbb Q(v),&v^4+pv^2+25p&=0,\\
 F_+&=\mathbb Q(w),&w^4-3pw^2+225p&=0,\\
 L&=F_-(t),&N&=KF_-=KF_+,\qquad M=KL.
 \end{aligned}
 \tag{PCQ.1}
\]
The fields $F_-$ and $F_+$ are the fixed fields of the kernels of
$\psi$ and $\epsilon_K\psi$, respectively, in the indicated
cyclotomic extensions.  The notation does not mean their value fields.

\begin{theorem}[AC-G156-class: two exact class inputs at the quarter residues]
\label{AC-G156-class}
The fields in (PCQ.1) have the class properties
\[
 \begin{gathered}
 \operatorname{Cl}(L)[p]=0,\qquad
 \operatorname{Cl}(F_-)[p]=0,\qquad
 \operatorname{Cl}(F_+)[p]=0,\\
 \operatorname{Cl}(M)[p]_{C_3\text{-nontrivial}}=0,\qquad
 \operatorname{Cl}(N)[p]_{b^2=-1}=0.
 \end{gathered}
 \tag{PCQ.2}
\]
Here $b$ generates the quartic factor with $\psi(b)=-i$;
$C_3$ is the cubic subgroup of $\operatorname{Gal}(M/K)$.
In particular the last term includes both $b$-eigenvalues $i,-i$.
For the five and three verified units used below, let $m_L,m_+$
be their respective indices in the full unit groups modulo torsion.
The exact integer relations established are
\[
             h_L=91m_L,\qquad h_{F_+}=m_+,
                     \qquad p\nmid m_Lm_+.
 \tag{PCQ.3}
\]
The assertions $h_L=91$, $h_{F_+}=1$, complete unit bases, and
triviality of the whole $p$-class groups of $M$ or $N$ are not used
or asserted.  The original finite labels remain $\chi^{-j}$;
these quartic characters occur only in residual splitting fields.
\end{theorem}

\begin{proof}
\proofstep{1. Actual fields, orientations and local discriminants.}
We retain the exact Gaussian-period identities already proved in
AC-G147-class.  If $g_p$ sends a primitive $p$th root to its square,
$v=\eta_0-\eta_2$ and $v'=\eta_1-\eta_3$, those identities are
\[
 v^2=-(p+13\sqrt p)/2,\quad
 vv'=-5\sqrt p,\quad
 g_p(v)=v'=-\frac{v^3+219v}{65},\quad g_p^2(v)=-v.
 \tag{PCQ.4}
\]
The positive sign of the quadratic Gauss sum and the complete
finite period multiplication certificate are retained with their
exact source hashes.  The present field verifier rechecks the
polynomial, square-root and automorphism identities; it does not
claim a second period enumeration.
Since $2$ has order $268$ modulo $p$, the displayed generator has
$\psi(g_p)=-i$.  This identifies $F_-$ with the cyclic quartic
field of $\psi$.
Its unique quadratic subfield is $\mathbb Q(\sqrt p)$, and it is
imaginary with two complex pairs.  Its conductor and discriminant are
$p$ and $p^3$, by the primitive character conductors $p,p,p$.

Take $w=\sqrt{-3}\,v$ in the common compositum with $K$.
Then $w^2=3(p+13\sqrt p)/2$, and its conjugate square uses the
minus sign.  Both are positive and have product $225p$.
The cyclic quartic field is therefore totally real, and is the
field of the primitive even character $\epsilon_K\psi$ of conductor
$3p$.  More precisely its kernel in $C_2\times C_4$ is the single
order-two subgroup generated by $(c_K,g_p^2)$.
Here $c_K$ is conjugation on $K$ and fixes $F_-$.
The actual quartic actions are
\[
 b(v)=-\frac{v^3+219v}{65},\qquad
 b(w)=\frac{w^3-657w}{195},\qquad b^2(v)=-v,
 \quad b^2(w)=-w.
 \tag{PCQ.5}
\]
The CRT representative $271$ fixes the conductor-three factor
and is $2$ modulo $p$, hence has $(\epsilon_K\psi)$-value $-i$.
The field $F_+$ has discriminant $(3p)^2p=9p^3$.
The two quartic polynomials are Eisenstein at $p$, which also
proves their degree directly.  All these character, polynomial
and inverse-action checks are literal in the saved field dictionary.

The unique quadratic subfield of the $S_3$ splitting field $K(t)$
is $K$, different from $\mathbb Q(\sqrt p)$.
Thus $K(t)\cap F_-=\mathbb Q$, and
\[
 \operatorname{Gal}(M/\mathbb Q)=S_3\times C_4,\qquad
 [L:\mathbb Q]=12,\quad [M:\mathbb Q]=24,\quad[N:\mathbb Q]=8.
 \tag{PCQ.6}
\]
The signatures, roots of unity, and absolute discriminants needed
for the class formula are
\[
 \begin{array}{c|c|c|c|c}
 A&[A:\mathbb Q]&(r_1,r_2)&w_A&D_A\\ \hline
 L&12&(0,6)&2&3^4p^{11}=43227497953464440249376865989\\
 F_+&4&(4,0)&2&9p^3=175185981.
 \end{array}
 \tag{PCQ.7}
\]
The real field has only two roots of unity.  The maximal abelian
subfield of $L$ is $F_-$: in $S_3\times C_4$, the commutator and
any reflection together generate the $S_3$ factor.  Any root of
unity in $L$ therefore belongs to $F_-$.  Since this quartic
abelian field has conductor $269$, a root of unity other than
$\pm1$ would give a cyclotomic subfield of degree at least $268$.
This is impossible.

For completeness, the discriminant of $L$ follows locally without
using a provisional PARI discriminant.  Only $3,p$ ramify in its
two factors.  At $p$, the quartic and cubic extensions are totally
ramified of coprime degrees four and three, giving a tame degree-twelve
extension and discriminant exponent eleven.  At $3$, $p$ is a cube
in $\mathbb Q_3$: write $a=-(1+3y)$ and solve
$y+3y^2+3y^3=-30$, whose derivative is one modulo three.
The other factor of $X^3-p$ is a ramified quadratic with discriminant
exponent one.  Tensoring with the unramified quartic algebra from
$F_-$ gives total exponent four.  At $2$, both factors are unramified.
The same local computation shows that $M$ has one $p$-place,
with absolute ramification index twelve and residue degree two.
The field $N$ has one $p$-place with indices four and two.

\proofstep{2. Exact units, with no fundamental-basis assumption.}
The candidate generator for $L$ is $v-t$.
The exact twelve-dimensional rational regular representation verifies
its polynomial and both inverse maps to $t,v$.
For every supplied unit, it verifies the printed monic integral
annihilating polynomial with constant coefficient $1$ or $-1$,
both by its characteristic polynomial and by direct Horner evaluation.
Thus each element and its inverse are integral.  The original
rational denominators, including $269$, are retained throughout.
There are five such units.  Their coordinate vectors below use
$t^av^b$, with index $4a+b$ and $0\leq a<3$, $0\leq b<4$.
Each numerator vector is divided by the indicated $d_i$.
\[
\begin{aligned}
 d_{1}&=13,\quad\mathbf n_{1}=(-2411,0,-10,0,\\
 &\hspace{3em}0,0,0,0,\\
 &\hspace{3em}0,0,0,0),\\
 d_{2}&=104910,\quad\mathbf n_{2}=(9275514085,-2089094619,349655615,-83450256,\\
 &\hspace{3em}4004017925,232897779,148287595,9446796,\\
 &\hspace{3em}-842873840,14050677,-31362895,539043),\\
 d_{3}&=104910,\quad\mathbf n_{3}=(9275514085,2089094619,349655615,83450256,\\
 &\hspace{3em}4004017925,-232897779,148287595,-9446796,\\
 &\hspace{3em}-842873840,-14050677,-31362895,-539043),\\
 d_{4}&=1,\quad\mathbf n_{4}=(270,0,0,0,\\
 &\hspace{3em}-739,0,0,0,\\
 &\hspace{3em}108,0,0,0),\\
 d_{5}&=104910,\quad\mathbf n_{5}=(84781846350,59960359047,349655615,249023253,\\
 &\hspace{3em}35885345130,-6799782807,148287595,-28245393,\\
 &\hspace{3em}-7593744915,-385436919,-31362895,-1599996),\\
\end{aligned}
\tag{PCQ.8}
\]
The complete annihilators are in
\path{all_units_exact_certificate.json}, whose SHA256 is
\begin{quote}\ttfamily\small
7f81a932eb0e6c15d19da301995be7c0c\\
c1d4d6b891d63cc2217aa36ff99109e.
\end{quote}
A separate four-dimensional regular representation verifies the
three units in $F_+$, in the basis $1,w,w^2,w^3$:
\[
 \begin{array}{c|r|rrrr}
 i&d_i&\multicolumn{4}{c}{\mathbf n_i}\\ \hline
 1&390&-16440&699&185&-7\\
 2&78&-26571&2901&37&-4\\
 3&39&837&0&-10&0.
 \end{array}
 \tag{PCQ.9}
\]
Their monic relations, exact checks and residue data are in
\path{quartic_plus/finite_certificate.json}, SHA256
\begin{quote}\ttfamily\small
b8b5940291e8ea7e67757daab9678b16\\
c1554a6c4391cf6d723dc544534d8296.
\end{quote}
Neither verifier uses a provisional class number or fundamental-unit
claim.  The later positive regulator determinants also directly
verify full real rank.

\proofstep{3. Finite residue maps prove the needed unit saturation.}
The exact verifier uses a rational prime $\ell\equiv1\pmod p$
with simple degree-one residue roots of the field equations.
Complete trial division verifies its primality; the prime divisors
of $\ell-1$ verify the displayed primitive root $g$.
Then $h=g^{(\ell-1)/p}$ has exact order $p$.
Evaluation of each integral unit at each root, followed by
$u\mapsto u^{(\ell-1)/p}$ and its exhaustive discrete logarithm
relative to $h$, gives a homomorphism from the unit group to
$\mathbb F_p$.  Roots are simple, so the maps really are local
residue maps; none of the original denominators is divisible by
these $\ell$.

For $L$, all twelve maps at $\ell=206593$ are recorded.
A selected five-row minor is
\[
 \begin{pmatrix}
 70&155&39&85&95\\
 199&95&65&85&39\\
 199&65&95&85&155\\
 70&39&155&85&65\\
 70&55&205&214&184
 \end{pmatrix},\qquad
 \det=-40364535504\equiv180\pmod p.
 \tag{PCQ.10}
\]
For $F_+$, the four maps at $\ell=38737$ have a three-row
minor of determinant $-11987025\equiv153\pmod p$.
The complete maps, unit residues, logarithms, row selections
and integer determinants are saved, not only their ranks.

Each residue homomorphism kills $p$th powers.  It also kills
the global roots of unity, whose order is prime to $p$.
The full ranks therefore show that the respective unit classes
are bases modulo $p$ of the full unit lattices, which have ranks
five and three by Dirichlet's theorem and (PCQ.7).
Their subgroups have finite indices $m_L,m_+$ and
\[
                          p\nmid m_Lm_+.
 \tag{PCQ.11}
\]
This is a $p$-saturation conclusion, not index one.

\proofstep{4. Exact zeta factors and finite Dirichlet values.}
Let $\rho$ be the standard two-dimensional representation of
$S_3$ attached to $X^3-p$.  Induction from a reflection gives
$1\oplus\rho$.  Hence the permutation representation for $L$
gives the exact Artin factorization
\[
 \zeta_L(s)=\zeta(s)L(s,\tau)|L(s,\psi)|^2
              L(s,\rho)L(s,\rho\tau)|L(s,\rho\psi)|^2
 \tag{PCQ.12}
\]
for real $s>1$, and by continuation at the residue at one.
Here squared absolute values abbreviate a pair of conjugate factors,
not a holomorphic absolute-value expression at complex $s$.
For the real quartic field one has
\[
       \zeta_{F_+}(s)=\zeta(s)L(s,\tau)
                                      |L(s,\epsilon_K\psi)|^2.
 \tag{PCQ.13}
\]
These identities use the actual primitive characters, including all
bad Euler factors.

The quadratic factor is evaluated by the finite formula
\[
 L(1,\tau)=-\frac1{\sqrt p}
       \sum_{a=1}^{p-1}\tau(a)\log(2\sin(\pi a/p)).
 \tag{PCQ.14}
\]
The elementary Gauss identity gives squared Gauss modulus $p$;
its real sign is fixed positively by the rigorous finite cosine sum
in the verifier, so the denominator is exactly $+\sqrt p$.
For the odd quartic character, the exact finite Bernoulli sum is
\[
 B_{1,\psi}=\frac1p\sum_{a=1}^{p-1}a\psi(a)=-5-i,
 \qquad |L(1,\psi)|^2=\frac{26\pi^2}{p}.
 \tag{PCQ.15}
\]
The first equality is checked by integer summation using
$\psi(a)=a^{67}$ in the four-root dictionary with $i\mapsto82$.
The second is the primitive odd Dirichlet functional equation at
one, with $L(0,\psi)=-B_{1,\psi}$; taking norms removes the phase.
No unknown period or class number is introduced.

For $q=3p$ and the primitive even character $\nu=\epsilon_K\psi$,
finite Fourier inversion gives
\[
 |L(1,\nu)|^2=\frac1q\left|
     \sum_{a=1}^{q-1}\overline{\nu(a)}
                         \log(2\sin(\pi a/q))\right|^2.
 \tag{PCQ.16}
\]
For example this follows by pairing the terms $a,-a$ in the
finite Fourier expansion of $-\log(1-re^{2\pi ia/q})$,
then taking $r\uparrow1$; every surviving sine is positive.
The primitive Gauss factor has squared norm $q$.
The conductor is precisely $3p$: the two factors have coprime
conductors three and $p$, and each has a nontrivial value on the
kernel removing its own prime.  The field verifier supplies the
explicit CRT witnesses.  The formulas (PCQ.14)--(PCQ.16) are the
same exact finite Dirichlet identities retained in the earlier
class certificates, with these newly checked conductors and parity.

\proofstep{5. The three actual rank-two conductors and Poisson formula.}
The quadratic field $K$ has class number one: its Minkowski bound
is $2\sqrt3/\pi<2$.  The cubic character over $K$ inducing $\rho$
is split at three by the local cube calculation in Step1 and is
primitive tame of order three at the inert prime $p$.
Twisting it by $1,\tau\circ\operatorname{Norm},
\psi\circ\operatorname{Norm}$ gives primitive finite Hecke
characters of conductors exactly $p\mathcal O_K$ in all three cases.
Indeed the residual norm map $\mathbb F_{p^2}^\times\to
\mathbb F_p^\times$ is onto, and the product inertia orders
are respectively three, six and twelve.  No cancellation can
remove the factor of order three.  All other primes are unramified.
The completed rank-two conductor is consequently always
\[
                           Q=|D_K|p^2=217083.
 \tag{PCQ.17}
\]
In particular no extra factor of $p$ is multiplied into $Q$ for
a twist already ramified at this same prime.

Write $a_n$ for the coefficients of $L(s,\rho)$.
At a good rational prime $q$, $a_q$ is the number of roots of
$X^3-p$ modulo $q$ minus one, and the local determinant is
$\chi_{-3}(q)$.  At $3$ the unramified line has eigenvalue one,
so $a_{3^k}=1$; at $p$ there are no inertia invariants, so
$a_{p^k}=0$ for $k>0$.  Multiplicativity and the degree-two
recurrence therefore give exactly the coefficients used by the
script.  Their eigenvalues are roots of unity, so
\[
                  |a_n|\leq d(n)\leq n,
       \qquad b_n=a_n\nu(n),\quad |b_n|\leq n.
 \tag{PCQ.18}
\]
Here $\nu=1,\tau,\psi$ in the three evaluations.

For clarity we replay the applied Poisson argument in this actual
prime-modulus case.  The lattice $\mathcal O_K$ has area
$\sqrt3/2$ and trace-dual scale $2/\sqrt{-3}$ for the real
Euclidean Fourier pairing.  The period lattice
$p\mathcal O_K$ has area $p^2\sqrt3/2$.
The finite Hecke character on its residue ring is trivial on
$\mathcal O_K^\times$ and primitive.  At zero frequency its sum
is zero.  Writing the perfect additive pairing as
$\operatorname{Tr}(-r\overline\beta/\sqrt{-3})$, substitution
$r\mapsto r/\overline\beta$ at a nonzero frequency gives the
factor $\theta(\overline\beta)^{-1}$ times one Gauss sum.
Conjugating the dual summation variable then gives
$\theta(\beta)^{-1}=\overline{\theta(\beta)}$, with the same
Gaussian norm.  At the only nonunit frequency, zero, the sum
is zero by nontriviality.
Finite Parseval gives Gauss modulus $p$.
The inverse different makes the additive pairing perfect; $3$ is
a unit modulo $p$.  This is the exact retained Gaussian-period
Poisson normalization, now over the field $\mathbb F_{p^2}$.

Put $a=2\pi/\sqrt Q$ and
$\Theta_\nu(y)=\sum_{n\geq1}b_n e^{-any}$.
Because $K$ has class number one, this is the corresponding
lattice sum divided by its six units.  The Gaussian Fourier transform
and the displayed lattice area give scalar $1/(py)$ and dual
exponent $-a|\beta|^2/y$.  Combining this scalar with Gauss modulus
$p$, and the conjugation of the dual variable just checked, gives
\[
 \Theta_\nu(y)=\frac{\varepsilon_\nu}{y}
             \overline{\Theta_\nu(1/y)},\qquad
                         |\varepsilon_\nu|=1.
 \tag{PCQ.19}
\]
The bars conjugate coefficients; the argument $y$ is real positive.
No root-number convention is guessed: the computed ball for
$T_\nu=\Theta_\nu(1)$ excludes zero, so
$\varepsilon_\nu=T_\nu/\overline{T_\nu}$ exactly.
For the two real coefficient sequences this is one.
Splitting the Mellin integral at one, using exponential decay at
both ends supplied by (PCQ.19), yields
\[
 L(1,\rho\nu)=\sum_{n\geq1}\frac{b_n e^{-an}}n
       +\varepsilon_\nu a
                         \sum_{n\geq1}\overline{b_n}E_1(an),
 \qquad E_1(x)=\int_x^\infty\frac{e^{-t}}t\,dt.
 \tag{PCQ.20}
\]
This is precisely the formula evaluated by the saved script.

\proofstep{6. Regulator normalization, outward errors and integer isolation.}
For $L$, take the six complex-place representatives
\[
 t=p^{1/3}\omega^j,\quad
 v=i\sqrt{(p+13\sqrt p)/2}\ \hbox{or}\
 i\sqrt{(p-13\sqrt p)/2},\qquad j=0,1,2.
 \tag{PCQ.21}
\]
The full weighted logarithm matrix has entries $2\log|u|$.
Delete its last row and take the absolute determinant to get $R_U$.
For $F_+$, use the four real roots
$\pm\sqrt{3(p+13\sqrt p)/2}$,
$\pm\sqrt{3(p-13\sqrt p)/2}$, with entries $\log|u|$,
and delete the last row.  Call its determinant $R_{U,+}$.
All norm-square balls exclude zero; the complete column sums
contain zero, as required by the unit norm relation.
The projected weighted-log convention gives exactly
\[
                      R_U=m_LR_L,\qquad
                      R_{U,+}=m_+R_{F_+}.
 \tag{PCQ.22}
\]
No factor for the deleted row or the complex places is omitted.
The regulator calculation uses $1500$ bits for $L$ and $512$ bits
for $F_+$.  Its recorded positive intervals are near
$82602737.7956$ and $444.992803468$, respectively; these
approximate decimals are not proof inputs.

For (PCQ.20), put $z=e^{-a}$ and $A_N=e^{-a(N+1)}$.
The complete omitted tails of the first sum, the second before
multiplication by $a$, and theta at one are at most
\[
 \frac{A_N}{1-z},\qquad
 \frac{A_N}{a(1-z)},\qquad
                 \frac{A_N((N+1)-Nz)}{(1-z)^2}.
 \tag{PCQ.23}
\]
These follow from (PCQ.18) and $E_1(x)\leq e^{-x}/x$.
All three cutoffs are $N=12000$ at $256$-bit precision.
The script adds the entire error bound to every applicable real
and imaginary component before division for the phase or any
subsequent product.  The finite Dirichlet sums have no infinite
tail.  All evaluations use outward Arb balls, including logarithm,
exponential, square root and $E_1$.

Let $Z_L,Z_+$ denote the residues at one given by (PCQ.12)--(PCQ.16)
and (PCQ.20).  The analytic class-number formula, with (PCQ.7), is
\[
 h_LR_L=\frac{2\sqrt{D_L}}{(2\pi)^6}Z_L,
 \qquad h_{F_+}R_{F_+}=\frac{\sqrt{D_{F_+}}}{8}Z_+.
 \tag{PCQ.24}
\]
Dividing only by the computed nonzero regulator determinants
therefore computes $r_L=h_L/m_L$ and $r_+=h_{F_+}/m_+$.
It does not assume a complete unit basis.

We include the exact denominator bound used for rational isolation.
For a number field $A$ of degree $n$, signature $(r_1,r_2)$
and absolute discriminant $D$, put
\[
 U_A=\left\lceil\frac{n!4^{r_2}(\lfloor\sqrt D\rfloor+1)}
                              {n^n3^{r_2}}\right\rceil,
 \qquad H_A=4^n\left\lceil\sqrt[3]{U_A^4}\right\rceil.
 \tag{PCQ.25}
\]
Minkowski and $\pi>3$ put an ideal of norm at most $U_A$
in every class.  The number of such ideals is bounded by
$U_A^{4/3}\zeta_A(4/3)$.
The Euler factors give $\zeta_A(s)\leq\zeta(s)^n$ for $s>1$,
and $\zeta(4/3)\leq1+\int_1^\infty x^{-4/3}\,dx=4$.
Consequently $h_A\leq H_A$.
If $r=h_A/m>1$, then $m<H_A$; an enclosure
$|r-c|<1/H_A$ for an integer $c$ forces $h_A=cm$.
For a selector one, the weaker $r>1/2$ suffices:
then $m<2H_A$ and $|r-1|<1/(2H_A)$ forces $h_A=m$.
Both conclusions follow by multiplying by $m$ and using that
an integer of absolute value less than one is zero.

The actual exact bounds and strict interval comparisons are
\[
 \begin{array}{c|r|r}
 A&U_A&H_A\\ \hline
 L&62758778279&4184265209389909016576\\
 F_+&1241&3414272
 \end{array}
 \tag{PCQ.26}
\]
and
\[
 \begin{aligned}
 r_L&>1,&|r_L-91|&<2.70\cdot10^{-57},
                  &1/H_L&>2.38\cdot10^{-22},\\
 r_+&>1/2,&|r_+-1|&<5.6\cdot10^{-149},
                  &1/(2H_{F_+})&>1.46\cdot10^{-7}.
 \end{aligned}
 \tag{PCQ.27}
\]
The saved tests compare the full balls and exact integers;
the displayed rounded bounds are only strict consequences.
Thus $h_L=91m_L$ and $h_{F_+}=m_+$.
Since $269\nmid91$ and (PCQ.11) holds, both $p$-class groups
vanish.  This proves the first and third assertions of (PCQ.2).
The integer selectors $91$ and $1$ enter only these final tests;
no provisional class or regulator assertion was used earlier.

\proofstep{7. Precisely the reflected and abelian class sectors.}
Extension of ideals from $F_-$ to $L$ followed by norm back is
multiplication by three.  As $3$ is prime to $p$ and
$\operatorname{Cl}(L)[p]=0$, this proves
$\operatorname{Cl}(F_-)[p]=0$.
Let $s\in S_3$ be the reflection whose fixed field in $M$
is $L$.  Norm to $L$ followed by extension acts on
$\operatorname{Cl}(M)[p]$ as $1+s$, and is zero.
Hence $s=-1$ on that module.  Its conjugates have isomorphic
fixed fields and also act by minus one.  Their products generate
$C_3$, which therefore acts trivially.  This proves precisely
the nontrivial cubic sector assertion in (PCQ.2).

Finally $\operatorname{Gal}(N/\mathbb Q)=\langle c_K\rangle
\times\langle b\rangle=C_2\times C_4$.
The fixed field of $c_K$ is $F_-$, and that of $c_Kb^2$
is $F_+$; the latter involution is actual complex conjugation.
Both norm-extension composites $1+c_K$ and $1+c_Kb^2$
vanish on $\operatorname{Cl}(N)[p]$ by the two proved quartic
class statements.  On the submodule where $b^2=-1$, these
operators are $1+c_K$ and $1-c_K$; adding them gives two.
Since $p$ is odd, that submodule is zero.  This proves the last
assertion of (PCQ.2) without assuming any of the other class
sectors vanish.
\end{proof}

\paragraph{Original coefficients and subsequent gates.}
For the actual finite labels $\chi^{-j}$, the residual Kummer
quartic factors at effective residues $66$ and $200$ are
$\psi^{-1}$ and $\psi$.
For $j=1,2$ the nontrivial cubic class input in $M$ applies;
for $j=0$ the odd quartic class input in $N$ is available.
This statement supplies class inputs only.  The actual local
unit images on the unique completions, both inverse quartic
projectors, original-$O$ lifting and the common-window comparison
must be proved separately.  No $\psi$ or $\epsilon_K\psi$
is inserted as an original finite host label.
The all-prime index bound, mixed realization $\mathcal R_{\rm cub}$
and H04.4 remain open.

\paragraph{Source and computation scope.}
The exact Gaussian periods and sign are the previously fully
verified AC-G147-class input; copies and hashes are retained.
The finite Fourier and theta argument is replayed above under
its exact new prime-modulus and parity hypotheses.
The analytic class-number formula and projected weighted-log
normalization retain Sutherland, MIT18.785 Lecture19 (2019),
pp.1 and4.  The primitive Dirichlet functional equation retains
DLMF25.15.3,25.15.5--6.
Minkowski, Dirichlet units, class extension/norm, Artin induction
and the abelian conductor--discriminant formula retain their
previously audited scopes; they are not claimed as new proofs
from implementation code.  The exact finite scripts use rational
regular representations and complete residue calculations.
The analytic scripts use documented python-flint0.9.0 outward
balls, retaining the Arb implementation as an input.
Candidate PARI extraction supplies rational elements and final
selectors only.  No GRH, \texttt{bnfcertify}, complete-unit or
exact full class-number assertion is used.
